\documentclass[11pt]{article}
\usepackage{mathblatt}


\begin{document}
\pfheftstil
\pfheftkopf{Grundwert $G$}{P10}
\pfrueckblick
\begin{pfaufg}{}{1.}{}
Ergänze die Tabelle Prozent | Bruch | Dezimalzahl.
\par\smallskip\noindent \begingroup\renewcommand{\arraystretch}{1.3}\begin{tabular}{|c|c|c|}\hline \textbf{Prozent} & \textbf{Bruch} & \textbf{Dezimalzahl} \\ \hline 1 \% & \rule{0pt}{3ex}\hspace{1.6cm} & \rule{0pt}{3ex}\hspace{1.6cm} \\ \hline 10 \% & \rule{0pt}{3ex}\hspace{1.6cm} & \rule{0pt}{3ex}\hspace{1.6cm} \\ \hline 20 \% & \rule{0pt}{3ex}\hspace{1.6cm} & \rule{0pt}{3ex}\hspace{1.6cm} \\ \hline 25 \% & \rule{0pt}{3ex}\hspace{1.6cm} & \rule{0pt}{3ex}\hspace{1.6cm} \\ \hline 50 \% & \rule{0pt}{3ex}\hspace{1.6cm} & \rule{0pt}{3ex}\hspace{1.6cm} \\ \hline 100 \% & \rule{0pt}{3ex}\hspace{1.6cm} & \rule{0pt}{3ex}\hspace{1.6cm} \\ \hline\end{tabular}\endgroup\par
\fusshilfe{\mbox{1:~0,01; 0,1; 0,2; 0,25; 0,5; 1}}
\end{pfaufg}
\pfstufekopf[sgrundwert]{Grundwert $G$}{in 2 der letzten 5 Prüfungen}
\pfgruppe{erst verstehen – nicht rechnen}
\begin{pfaufg}{}{2.}{}
Ein Pullover kostete $60\,€$. Jetzt kostet er $45\,€$. Welcher Preis ist hier das Ganze?
\fusshilfe{\mbox{2:~der alte Preis, $60\,€$}}
\end{pfaufg}
\begin{pfaufg}{}{3.}{}
30 \% sind 21 €. Ist das Ganze größer oder kleiner als 21 €? Was kann passen? Kreuze an, rechne nicht.
\pfkreuzzeile{6,30 €}\pfkreuzzeile{21 €}\pfkreuzzeile{70 €}\pfkreuzzeile{210 €}
\fusshilfe{\mbox{3:~größer; 70 € (21 € : 0,3)}}
\end{pfaufg}
\par\addvspace{4pt}{\small\uebersichtskasten{\textbf{Merke:} Grundwert = das Ganze = 100 \%. Ist p kleiner als 100 \%, ist das Ganze größer als der Teil.}}\par
\pfgruppe{}
\begin{pfaufg}{}{4.}{}
Der graue Abschnitt ist $20\,\%$ des Ganzen. Trage den Abschnitt so oft ab, bis du bei $100\,\%$ bist. Markiere das Ganze.
\par\smallskip\noindent \streifen[0]{20}{$0\,\%$}{}\par
\par\noindent \leerfeld[Abschnitte] \par
\fusshilfe{\mbox{4:~$5$ Abschnitte bis $100\,\%$}}
\end{pfaufg}
\begin{pfaufg}{}{5.}{}
Ein Händler hat $24$ kg Äpfel verkauft. Das sind $10\,\%$ seiner Äpfel. Wie viel kg Äpfel hatte er?
\par\noindent \leerfeld[kg] \par
\fusshilfe{\mbox{5:~$240$ kg}}
\end{pfaufg}
\begin{pfaufg}{}{6.}{}
$7\,\%$ eines Betrags sind $56\,€$. Rechne in der Tabelle: erst $1\,\%$, dann das Ganze ($100\,\%$).
\par\smallskip\noindent \begingroup\renewcommand{\arraystretch}{1.35}\begin{tabular}[t]{|r|r|}\hline \textbf{Prozent} & \textbf{€} \\ \hline $7\,\%$ & $56$\,€ \\ \hline $1\,\%$ & \leerzelle \\ \hline $100\,\%$ & \leerzelle \\ \hline\end{tabular}\endgroup\hspace{2mm}{\small\color{mbgrau}\begin{tabular}[t]{@{}l@{}}\rule{0pt}{2.6ex}\\ $\downarrow : 7$\\ $\downarrow \cdot 100$\end{tabular}}\par
\fusshilfe{\mbox{6:~$800\,€$}}
\end{pfaufg}
\begin{pfaufg}{}{7.}{}
$36\,\%$ einer Menge sind $81$ kg. Wie viel kg ist die ganze Menge?
\par\noindent \leerfeld[kg] \par
\fusshilfe{\mbox{7:~$225$ kg}}
\end{pfaufg}
\pfgruppe{}
\pfab{$G = W : p$}
\begin{pfaufg}{P10 ’23}{8.}{}
Mia gewinnt bei einem Wettbewerb Geld. Sie gibt 25 \% des Gewinns an ihre Eltern ab; das sind 200 €. Gib an, wie hoch der gesamte Gewinn war.
\rechenplatz{1}
\fusshilfe{\mbox{8:~800 €}, \mbox{Tipp: $G$ = 200 € : 0,25}}
\end{pfaufg}
\par\addvspace{4pt}\Needspace*{0.3\textheight}\begin{list}{}{\setlength{\leftmargin}{8mm}\setlength{\topsep}{0pt}}\item[]
\noindent{\footnotesize\color{mbgrau}gleiche Art \textperiodcentered{} 2$\times$ geprüft – kannst du überspringen}\par
\begin{pfaufg}{P10 ’25}{9.}{}
Ein T-Shirt ist um 6 € billiger geworden. Dieser Nachlass entspricht 20 \% des ursprünglichen Preises. Gib den ursprünglichen Preis an.\par\vspace{7mm}\noindent{\color{black!45}\rule{\linewidth}{0.4pt}}
\fusshilfe{\mbox{9:~30 €}, \mbox{Tipp: $G$ = 6 € : 0,2}}
\end{pfaufg}
\end{list}
\pfgruppe{}
\begin{pfaufg}{BW ’23}{10.}{}
Am „Black Friday“ kostet ein Tablet 360,00 €. Der Preis wurde auf 80 \% gesenkt. Was kostete das Tablet vorher?
\fusshilfe{\mbox{10:~450 € (360 € sind 80 \%)}}
\end{pfaufg}
\pfgruppe{}
\begin{pfaufg}{nach P10 ’19}{11.}{}
In einem Gefäß liegen 4 schwarze Kugeln. Sie sind $\tfrac{2}{3}$ aller Kugeln im Gefäß. Berechne, wie viele Kugeln insgesamt im Gefäß liegen.
\rechenplatz{1}
\fusshilfe{\mbox{11:~G = 4 : $\tfrac{2}{3}$ = 6 Kugeln ($\tfrac{1}{3}$ sind 2 Kugeln)}}
\end{pfaufg}
\pfgruppe{}
\begin{pfaufg}{nach P10 ’17}{12.}{}
In einem Becher steht der Quark 8 cm hoch; das sind 93 \% der Becherhöhe. Berechne die Höhe des Bechers.
\rechenplatz{1}
\fusshilfe{\mbox{12:~G = 8 cm : 0,93 $\approx$ 8,60 cm}}
\end{pfaufg}
\pfstufekopf[sschluss]{Zum Schluss}{zwei ältere Aufgaben, damit nichts verloren geht}
\begin{pfaufg}{P10 ’18}{13.}{}
Für ein Fest backt Eva 16 Pfannkuchen: 14 füllt sie mit Marmelade, 2 mit Senf. Gib an, wie viel Prozent der Pfannkuchen eine Senffüllung haben.
\rechenplatz{1}
\fusshilfe{\mbox{13:~12,5 \%}}
\end{pfaufg}
\begin{pfaufg}{P10 ’21}{14.}{}
Ein Fahrrad kostet 550 €. Wer bar bezahlt, bekommt 20 \% Rabatt. Kreuze an, wie viel Euro man dabei spart:
\pfkreuzzeile{100 €}\pfkreuzzeile{440 €}\pfkreuzzeile{110 €}\pfkreuzzeile{660 €}
\fusshilfe{\mbox{14:~110 €}}
\end{pfaufg}
\end{document}